\documentclass[10pt,a4paper]{amsart}
\usepackage[margin=19mm]{geometry}
\usepackage{amsmath,amssymb,mathtools}
\usepackage[scr=rsfs,cal=euler]{mathalfa}
\usepackage{xfrac}
\usepackage[unicode,hidelinks,bookmarks=false]{hyperref}
\newcommand{\ie}{i.e.\ }
\newcommand{\cf}{cf.\ }
\newcommand{\resp}{respectively\ }
\newcommand{\Subsection}{\subsection}
\providecommand{\nobreakdash}{\nobreak\hbox{-}}
\setlength{\emergencystretch}{2em}
\multlinegap0pt
\usepackage{bm}
\usepackage{bbm}
\usepackage{dsfont}
\usepackage{xspace}
\usepackage{cancel}
\usepackage{mathtools}
\usepackage{amsthm}
\makeatletter
\@ifundefined{theorem}{\newtheorem{theorem}{Theorem}}{}
\@ifundefined{proposition}{\newtheorem{proposition}[theorem]{Proposition}}{}
\makeatother
\title[Existence and spectral stability of small-amplitude periodic]{Existence and spectral stability of small-amplitude periodic waves for the 2D nonlinear focusing Schrodinger equation}
\author{Fabio Natali}
\date{Source: arXiv:2503.13201v3 (CC BY 4.0)}
\begin{document}
\maketitle

\noindent\emph{Source and licence.} Existence and spectral stability of small-amplitude periodic waves for the 2D nonlinear focusing Schrodinger equation. Version arXiv:2503.13201v3 (CC BY 4.0), distributed under the Creative Commons Attribution 4.0 International licence, as recorded in the arXiv licence metadata for this exact version and shown on the abstract page https://arxiv.org/abs/2503.13201v3. Full rights notice: licenses/arxiv-2503.13201-CC-BY-4.0.txt.\par\smallskip

\noindent\textit{Editorial scope.} Counted unit: the Proposition whose source label is \texttt{cor-existence} in the article (source0.tex lines 542--546, source label \texttt{cor-existence}), delivered together with the article's own complete original proof of it (source0.tex lines 548--593, one \texttt{proof} environment of 46 lines). No mathematics is written, repaired or re-derived: statement and proof are the article's own text line for line. Required context retained uncounted: infB (lines 511--515); ode-wave (lines 523--525); minlem (lines 529--536). Every definition, display and label of those lines is retained unchanged. Omitted from this package and not redistributed: 1--510, 516--522, 526--528, 537--541, 547--547, 594--1062. Nothing inside a retained excerpt is omitted or altered: every retained body is delivered exactly as the article's lines are written, so no omission receipt of any kind accompanies this package. Citations: the retained excerpts carry no citation command at all, so no bibliography entry of the article is transcribed and no reference is redistributed. Reference adaptation: none was needed and none was made; every reference inside the delivered unit resolves inside it, so no unresolved reference or citation is produced. Printed slips kept literal: no printed word, symbol, hypothesis or number of the counted statement or of its proof was altered. Layout and class: the article's own class is \texttt{amsart} with packages including graphicx, amssymb, amsmath, amsthm, color, hyperref, mathabx, subfigure, lineno, verbatim, relsize; the wrapper is the lane's pinned \texttt{amsart} editorial wrapper at 10pt on A4, which restates the theorem-like environments the retained text needs and adds the article's own macro definitions that the retained text needs (none), with extra wrapper packages (bm, bbm, dsfont, xspace, cancel, mathtools). The delivered numbering is the unit's own. Assets: no retained span contains a figure environment, a graphics-inclusion command, a PDF-page inclusion, a table or a TikZ picture; no image file is copied into this package, the package declares no asset dependency and no image file is redistributed with it. Licence: version arXiv:2503.13201v3 of this article is distributed under the Creative Commons Attribution 4.0 International licence, as recorded in the arXiv licence metadata for this exact version and shown on the abstract page https://arxiv.org/abs/2503.13201v3. A full rights notice is delivered at licenses/arxiv-2503.13201-CC-BY-4.0.txt. Characters: every retained excerpt is pure ASCII, so the delivered engine, XeLaTeX with its default Latin Modern text font, reports no missing character.
			\begin{equation}
				\label{infB}
				q = \inf_{u\in Y_\lambda} B_c(u), \quad
				B_c(u) = \frac{1}{2}\int_{\mathbb{T}\times\mathbb{T}}|\nabla u|^2+c|u|^2 dxdy,
			\end{equation}

			\begin{equation}\label{ode-wave}
			-\Delta \Psi+c\Psi -|\Psi|\Psi=0.	
				\end{equation}

\begin{proposition}
\label{minlem}
There exists a ground state
of the constrained minimization problem \eqref{infB}. In other words, there exists $\Psi\in Y_{\lambda}$ satisfying
\begin{equation}\label{minBfunc}
B_c(\Psi)=\inf_{u\in Y_{\lambda}}B_c(u).
\end{equation}
\end{proposition}

\begin{proposition}
\label{cor-existence}
There exists a solution to the periodic boundary-value problem
(\ref{ode-wave}) with the minimizer profile $\Psi$ obtained in Proposition $\ref{minlem}$. In addition, $\Psi$ can be consider as $\Psi(x,y)=e^{i\theta}\varphi(x,y)$, where $\theta\in\mathbb{R}$ and $\varphi>0$.
\end{proposition}

\begin{proof}
By Lagrange's multiplier theorem, the constrained minimizer $\Phi \in Y_{\lambda}$ in Proposition \ref{minlem} satisfies
the stationary equation
\begin{equation}
\label{lagrange}
-\Delta \Psi+c\Psi=C_1|\Psi|\Psi
\end{equation}
for some constant $C_1$. A standard scaling argument enables us to consider $C_1=1$.\\
\indent Next, consider $\Psi=\psi_1+i\psi_2$ and define $\Phi=|\psi_1|+i|\psi_2|$. We see that $|\Psi|=|\Phi|$ and $|\nabla\Psi|=|\nabla\Phi|$, and both two facts imply that $\Phi$ is a also a minimizer of the problem $(\ref{infB})$. We see that $(\ref{lagrange})$ is also satisfied for $\Phi$ in place of $\Psi$. Separating equation $(\ref{lagrange})$ into real and imaginary parts, we obtain that $|\psi_1|$ and $|\psi_2|$ satisfies the following equations

\begin{equation}\label{psi11}
-\Delta |\psi_1|+c|\psi_1|=|\Psi||\psi_1|,
\end{equation}
and
\begin{equation}\label{psi22}
-\Delta |\psi_2|+c|\psi_2|=|\Psi||\psi_2|.
\end{equation}
Equalities $(\ref{psi11})$ and $(\ref{psi22})$ give us, since $c>0$, that $|\psi_1|$ and $|\psi_2|$ are elements in $H_{\rm per,s}^2\hookrightarrow C_{\rm per,s}$. Here $C_{\rm per,s}$ indicates the space of periodic continuous function $f$ satisfying $f(-x,y)=f(x,y)$ and $f(x,-y)=f(x,y)$ for all $(x,y)\in\mathbb{T}\times\mathbb{T}$. Using the strong maximum principle, we obtain, since $|\psi_1|,|\psi_2|\geq0$, that $|\psi_1|,|\psi_2|>0$ over $\mathbb{T}\times\mathbb{T}$. This last fact establishes that $\psi_1$ and $\psi_2$ do not change their sign over $\mathbb{T}\times\mathbb{T}$. We can assume, without loss of generality that both $\psi_1$ and $\psi_2$ are positive over $\mathbb{T}\times\mathbb{T}$.\\
\indent We claim that $\psi_1=\alpha\psi_2$ for some $\alpha\in\mathbb{R}$. Suppose that such an $\alpha\in\mathbb{R}$ does not exist. Define the real-valued function $z=\psi_1-\psi_2$. By linearity and the fact that $\psi_1$ and $\psi_2$ also solve equation $(\ref{psi11})$ and $(\ref{psi22})$, respectively, we obtain that function $z$ solves the equation
\begin{equation}\label{z}
-\Delta z+cz=|\Psi|z.
\end{equation}
It is important to mention that we can suppose, without loss of generality, that $\psi_1 \perp \psi_2$, since ${\psi_1, \psi_2}$ is linearly independent and both functions satisfy the linear equations $(\ref{psi11})$ and $(\ref{psi22})$, respectively. Since $(z, \psi_1)_{L_{\rm per}^2} = ||\psi_1||_{L_{\rm per}^2}^2 > 0$ and $(z, \psi_2)_{L_{\rm per}^2} = -||\psi_2||_{L_{\rm per}^2}^2 < 0$, we conclude that $z$ changes its sign over $\mathbb{T} \times \mathbb{T}$. Suppose, without loss of generality, that $z>0$. Since $\psi_1$ and $\psi_2$ are strictly positive functions over $\mathbb{T}\times\mathbb{T}$, we obtain, by continuity, that both inner products $(z,\psi_1)_{L_{\rm per}^2}$ and $(z,\psi_2)_{L_{\rm per}^2}$ are positive. This leads a contradiction, since $(z, \psi_2)_{L_{\rm per}^2} = -||\psi_2||_{L_{\rm per}^2}^2<0$. \\
\indent On the other hand, let us consider the (real)  minimization problem
\begin{equation}
				\label{infB1}
				b = \inf_{u\in W_\tau} D(u), \quad
				D(u) = \frac{1}{2}\int_{\mathbb{T}\times\mathbb{T}}|\nabla u|^2-|\Psi|u^2 dxdy,
			\end{equation}
			where $W_{\tau}=\left\{u\in H_{\rm per,s}^1;\ \int_{\mathbb{T}\times\mathbb{T}} u^2 dxdy = \tau \right\}$. A standard compactness argument, shows that the minimum of the problem $(\ref{infB1})$ is attained in a periodic real function $v$. Moreover, from Lagrange's multiplier theorem, we guarantee the existence of $d\in\mathbb{R}$ such that
\begin{equation}\label{lambda}
-\Delta v+dv=|\Psi|v.
\end{equation}

 		
By defining $w=|v|$ and doing a similar procedure as done for $\Phi$ that $w$ is a also a solution of the minimization problem $(\ref{infB1})$ and it satisfies the following linear equation
\begin{equation}\label{w}
-\Delta w+d w=|\Psi|w.
\end{equation}
In addition, integrating equation $(\ref{w})$ over $\mathbb{T}\times\mathbb{T}$, we obtain automatically that $d>0$. Thus, from the facts that $d>0$ and $w\geq0$, and by applying the strong maximum principle, we obtain that $w>0$, meaning that $v$ does not change its sign. Multiplying $(\ref{lambda})$ by $\psi_1$ and integrating over $\mathbb{T}\times\mathbb{T}$, we obtain after integrating by parts 
\begin{equation}\label{d-c}
		(d-c)\int_{\mathbb{T}\times\mathbb{T}}v\psi_1dxdy=0.
\end{equation}
Since both $\psi_1$ and $v$ are positive periodic functions, we deduce that $d=c>0$. This fact enables us to conclude that the all minimizers of the problem $(\ref{infB1})$ can be considered positive and they solve equation $(\ref{z})$.\\
\indent Next, consider $h=\frac{\tau z}{||z||_{L_{per}^2}}$. Multiplying equation $(\ref{z})$ by $\frac{z}{||z||_{L_{per}^2}^2}$ and integrating the result over $\mathbb{T}\times\mathbb{T}$, we obtain $D(h)=-\frac{c\tau}{2}=b$. This implies that $h$ is also a minimizer of problem $(\ref{infB1})$. Using a similar analysis as above, we can conclude that $|h|$ is also a minimizer, and by the strong maximum principle, we obtain that $|h| > 0$, meaning that $h$ does not change its sign over $\mathbb{T} \times \mathbb{T}$, nor does $z$. This leads to a contradiction because we claimed that $z=\psi_1-\psi_2$ changes its sign. Thus, we have $\psi_1 = \alpha \psi_2$ with $\alpha > 0$. Finally, since $\Psi = \psi_1 + i\psi_2 = (\alpha + i)\psi_2$, we can consider $\theta \in \mathbb{R}$ such that $\Psi = e^{i\theta}\varphi$ with $\varphi > 0$, as desired. This completes the proof of the proposition.
\end{proof}

\end{document}
